% Independent AI image-based transcription; QA evidence only.
% Source: extra/SGA-archive/sga7-1.pdf, physical pages 2--4.
% NOT included in the deliverable PDF. Paragraphs and mathematics preserved;
% line-break hyphenation joined, accents normalized in UTF-8.
% Page 2 has no visible original Roman folio; pages 3,4 read VI,VII.
\section*{INTRODUCTION}

% Physical PDF page 2; original introduction first page.
Soient $S$ un schéma et $f:X\longrightarrow S$ un morphisme de schémas.
Si $f$ est propre et lisse, et que le nombre premier $\ell$ est inversible
sur $S$, les groupes de cohomologie $\ell$-adique
$H^i(X_{\bar s},\mathbb Z_\ell)$ des fibres géométriques de $f$ forment
un système local $\ell$-adique sur $S$. Pour $f$ seulement supposé propre,
ces groupes sont les fibres d'un faisceau $\ell$-adique
$R^if_*\mathbb Z_\ell$ sur $S$. La \underline{théorie des cycles évanescents}
met en relation la ramification de ce faisceau sur $S$ et les singularités de $f$.

Nous ne considérerons que le cas où $S$ est un trait hensélien
(= spectre d'un anneau de valuation discrète hensélien). C'est en pratique
le cas essentiel. Je renvoie à (I 2.1) pour une description heuristique
de la théorie. Soient $S=\operatorname{Spec}(V)$, $k(\bar\eta)$ une clôture
algébrique du corps des fractions $k(\eta)$ de $V$ et $k(\bar s)$ la clôture
algébrique correspondante du corps résiduel $k(s)$
($k(\bar s)$ est le corps résiduel du normalisé $\bar V$ de $V$ dans
$k(\bar\eta)$). Dans le cas particulier où $X$ est propre et plat sur $S$,
de dimension relative $n$, et où $f$ ne présente qu'un point de non lissité
$x\in X_s$, on définit des $\operatorname{Gal}(\bar\eta/\eta)$-modules
$\phi^i$, de nature purement locale au voisinage de $x$, nuls pour
$i\notin[0,n]$ (I 4.2), et on construit une suite exacte longue de
$\operatorname{Gal}(\bar\eta/\eta)$-modules
\[
\cdots\longrightarrow H^i(X_{\bar s},\mathbb Z_\ell)
\xrightarrow{\mathrm{sp}}H^i(X_{\bar\eta},\mathbb Z_\ell)
\longrightarrow\phi^i\longrightarrow
H^{i+1}(X_{\bar s},\mathbb Z_\ell)\longrightarrow\cdots
\]
(le $\operatorname{Gal}(\bar s/s)$-module $H^i(X_{\bar s},\mathbb Z_\ell)$
est regardé comme un $\operatorname{Gal}(\bar\eta/\eta)$-module avec action
triviale de l'inertie ; $\mathrm{sp}$ est la flèche de spécialisation).
On donne aussi des critères, valables pour l'instant seulement en
caractéristique $0$, pour que les $\phi^i$ pour $i$ petits soient nuls
(par exemple, si $X$ est de plus localement d'intersection complète,
$\phi^i=0$ pour $i\ne n$ (I. 4.5)).

% Physical PDF page 3; printed Roman folio VI.
Le \underline{théorème de monodromie} affirme que l'action du sous-groupe
d'inertie $I$ de $\operatorname{Gal}(\bar\eta/\eta)$ est quasi-unipotente :
pour $T\in I$, il existe des entiers $N>0$, $M>0$ tels que
l'endomorphisme $(T^M-1)^N$ de $H^i(X_{\bar\eta},\mathbb Z_\ell)$ soit nul.
On donne de ce théorème deux démonstrations. La première (I.1.2), de
nature arithmétique, s'applique dès que les groupes de cohomologie
considérés ont des propriétés de finitude raisonnables. Le point clef
est que, lorsque $k(s)$ est de type fini sur son sous-corps premier,
l'action de $I$ est quasi-unipotente pour toute représentation
$\ell$-adique de $\operatorname{Gal}(\bar\eta/\eta)$. La seconde
démonstration, plus géométrique, requiert la pureté et la résolution
des singularités : elle n'est pour l'instant valable qu'en caractéristique
$0$ ou pour un $H^1$. Elle apporte de précieuses informations sur
l'exposant de nilpotence $N$ ($N\leq i+1$ pour un $H^i$). Ainsi qu'on le
verra ultérieurement, elle se prête bien, sur $\mathbb C$, à la comparaison
avec la théorie transcendante.

Ces résultats, appliqués au $H^1$ des variétés abéliennes, i.e. à leur
module de Tate, permettent d'étudier la réduction mod $p$ de celles-ci.
On donne ainsi deux démonstrations du théorème de réduction stable
des variétés abéliennes selon lequel, après ramification, la fibre
spéciale connexe du modèle de Néron est extension d'une variété abélienne
par un tore. La première (I. 6) reprend la méthode de la démonstration
arithmétique du théorème de monodromie. La seconde (IX 3.6) s'appuie sur
une analyse beaucoup plus fine du modèle de Néron, et de ses propriétés
de polarisation.

Les exposés I à V de Grothendieck n'ont pas été rédigés. Ils ont été
résumés dans un exposé I. Les résultats énoncés y sont démontrés de
façon succincte, mais essentiellement complète.

L'exposé II applique la méthode des pinceaux de Lefschetz (I 5.3) à
l'étude du groupe fondamental. Pour $S$ une surface sur un corps
algébriquement clos $k$, d'exposant caractéristique $p$, on montre que
le groupe profini $\pi_1^{(p)}(S,s)$, complété en dehors de $p$ du groupe
fondamental, est de pro-(p)-présentation finie.

Comme expliqué plus haut, les exposés III à V n'existent pas.

% Physical PDF page 4; printed Roman folio VII.
L'exposé VI contient, avec quelques compléments, la théorie des
déformations de Schlessinger. On y prend soin de tenir compte des
automorphismes infinitésimaux des objets qu'on classifie, et de ne pas
passer trop brutalement au foncteur des classes d'isomorphie d'objets.
On y donne aussi une nouvelle construction des
$\underline{\operatorname{Ext}}^i(L_{X/S},-)$
($L_{X/S}$ complexe cotangent relatif de $X/S$). Cet exposé sert dans
le reste du séminaire surtout via l'étude qui y est faite des
déformations des singularités quadratiques ordinaires.

Les exposés VII et VIII sont consacrés à la théorie des biextensions.
Cette théorie joue un rôle essentiel dans l'exposé IX, pour exprimer
ce qu'il advient d'une polarisation quand on passe d'une variété abélienne
à un modèle de Néron de celle-ci. Cet exposé IX contient la démonstration
du théorème de réduction stable des variétés abéliennes, et diverses applications.

La suite de ce séminaire : SGA 7 II, par P. Deligne et N. Katz, paraîtra
ultérieurement.

\begin{flushright}
Bures sur Yvette, mai 1972,\\[1ex]
P. DELIGNE
\end{flushright}
